% !TEX root = ../../../main.tex
% Block card: catalogue.gramian_gasf -- "Gramian GASF" (Part I, the blocks chapter).
\cardfield{What it does}
Turns a short stretch of one channel into a square picture in which every pixel compares
two moments of the stretch with each other.

\begin{plainterms}[Gramian images and recurrence plots]
Picture the signal's lowest value as ``pointing right'' and its highest as ``pointing
left'', with every value in between a direction part of the way round a half circle. Each
moment of the stretch becomes such a direction. Now make a table with one row and one
column per moment, and in each cell write how two moments' directions combine. That table,
drawn in colour, is the image. The summation field (this block) adds the two directions; the
difference field (\blockref{catalogue.gramian_gadf}) subtracts them. The recurrence plot
(\blockref{catalogue.gramian_recurrence}) fills the same kind of table with how far apart two
short snatches of the signal are. The pictures let an image-recognition network look at a
signal \citep{wang2015imaging}.
\end{plainterms}

\cardfield{When to use it}
To look at the texture of one short stretch, or to feed an image-based model. For images of
many windows at once, as the trained networks expect, use \blockref{catalogue.window_images}.

\cardfield{What goes in, what comes out}
In goes a span of one channel, at most 5{,}000 samples long. Out comes a square image with
one row and one column per sample: a 600-sample span gives a $600 \times 600$ image.

\cardfield{Settings that matter}
None: the block has no settings. The span you choose is everything.

\cardfield{Watch out for}
The picture is stretched between the span's own lowest and highest values, so one big spike
sets the scale for everything else. A span that is perfectly flat, or that has a single
missing sample, gives an image with no values at all. A span longer than 5{,}000 samples is
refused.
